\section{The Project at a Glance}

\subsection{Goal}

We observe an $n\times d$ matrix $\underline{X}=(x_{ih})$ of \hl{ordinal data}:
each entry is a category $x_{ih}\in\{1,\ldots,m\}$ (for instance a rating
on a scale from 1 to $m$; rows are individuals, columns are items).

\hl{Co-clustering} simultaneously
\begin{itemize}
  \item groups the $n$ rows into $K$ row clusters,
  \item groups the $d$ columns into $L$ column clusters,
\end{itemize}
so that the matrix is summarised by $K\times L$ homogeneous \hl{blocks}.

\paragraph{What we have to deliver.}
An R package whose main function takes the data, $K$, $L$ and the number
of random initialisations (20 by default) and returns
the membership probabilities of rows and columns, the hard assignments,
the ICL value and the estimated parameters.

\paragraph{The recipe.}
\begin{center}
\begin{tabular}{ll}
\toprule
Ingredient & Choice \\
\midrule
Probabilistic model & Latent Block Model (LBM) \\
Distribution inside a block & CUB (ordinal) \\
Inference & SEM-Gibbs algorithm \\
Choice of $(K,L)$ & ICL criterion \\
Robustness to local optima & several random initialisations \\
\bottomrule
\end{tabular}
\end{center}

\paragraph{Why not simply EM?}
For the LBM the E step requires
$P(z,w\mid x)$, which does not factorise (rows and columns are
coupled). It is a sum over $K^nL^d$ configurations: intractable.
SEM-Gibbs replaces this sum by \hl{simulation}: we sample $(z,w)$ with
Gibbs sampling instead of computing expectations.

\section{The CUB Distribution}

\subsection{Definition}

For one ordinal variable $R\in\{1,\ldots,m\}$, the CUB model
(Combination of a Uniform and a shifted Binomial) is the mixture
\[
\boxed{
P(R=r)
=
\pi\binom{m-1}{r-1}\xi^{\,m-r}(1-\xi)^{\,r-1}
+(1-\pi)\frac1m,
\qquad r=1,\ldots,m.
}
\]

\paragraph{Interpretation.}
A respondent either
\begin{itemize}
  \item answers \hl{with feeling} (probability $\pi$): the answer follows a
  shifted binomial $1+\Bin(m-1,1-\xi)$, or
  \item answers \hl{with uncertainty} (probability $1-\pi$): the answer is
  uniform over the $m$ categories.
\end{itemize}
So $\pi$ is the weight of the feeling component and $\xi$ drives the
location of the answers: the mean of the binomial part is
$1+(m-1)(1-\xi)$, hence a \emph{large} $\xi$ means \emph{low} categories.

\paragraph{Check that it is a distribution.}
The binomial terms sum to one, $\sum_r\binom{m-1}{r-1}\xi^{m-r}(1-\xi)^{r-1}=1$,
and so do the uniform terms. Hence
$\sum_rP(R=r)=\pi+(1-\pi)=1$.

\subsection{Latent variable representation}

Introduce $V\sim\text{Bernoulli}(\pi)$ with
\[
P(R=r\mid V=1)=b(r;\xi)=\binom{m-1}{r-1}\xi^{m-r}(1-\xi)^{r-1},
\qquad
P(R=r\mid V=0)=\frac1m.
\]
This turns CUB into a two-component mixture and allows an EM algorithm.

\subsection{Estimation by EM (one CUB, counts $c_r$)}

Let $c_r$ be the number of observations equal to $r$, and $N=\sum_rc_r$.

\paragraph{E step.} Posterior probability of ``feeling'' given $R=r$:
\[
\boxed{
\nu_r
=
\frac{\pi\,b(r;\xi)}{\pi\,b(r;\xi)+(1-\pi)/m}.
}
\]

\paragraph{M step.} The expected complete log-likelihood is
\[
\sum_rc_r\Big[\nu_r\big(\log\pi+\log b(r;\xi)\big)
+(1-\nu_r)\big(\log(1-\pi)-\log m\big)\Big].
\]
Maximising in $\pi$:
\[
\boxed{\pi=\frac{\sum_rc_r\nu_r}{N}.}
\]
In $\xi$, only $\sum_rc_r\nu_r\big[(m-r)\log\xi+(r-1)\log(1-\xi)\big]$
matters. Setting the derivative to zero and using
$(m-r)+(r-1)=m-1$:
\[
\boxed{\xi=\frac{\sum_rc_r\nu_r(m-r)}{(m-1)\sum_rc_r\nu_r}.}
\]

\paragraph{Why counts?}
Only the $m$ numbers $c_r$ matter, not the individual data. This is the key
efficiency trick of the code: for each block $(k,\ell)$ we tabulate the
$m$ counts and run this tiny EM on them.

\paragraph{Pitfalls.}
\begin{itemize}
  \item If $\pi\to0$, $\xi$ is not identifiable (the data are uniform); the
  code clamps $\pi$ and $\xi$ to $[10^{-6},1-10^{-6}]$.
  \item If a block is empty ($N=0$), the parameters are kept unchanged.
\end{itemize}
